% Source problem: public/notes/物工/問題/2017ap_exam.pdf, pages 11–12.
\documentclass[11pt,a4paper]{article}
\usepackage[margin=24mm]{geometry}
\usepackage{fontspec}
\usepackage{xeCJK}
\setCJKmainfont{HaranoAjiMincho-Regular.otf}
\usepackage{amsmath,amssymb}
\usepackage{tikz}
\usetikzlibrary{arrows.meta}
\setlength{\parindent}{1em}
\setlength{\parskip}{0.35em}
\title{物理工学専攻 院試解答集\\2017年 物理学II 第4問}
\author{政岡凜太郎}
\date{}
\begin{document}
\maketitle

\section*{[1.1] 電子の伝導率テンソル}
$x=\omega_c\tau$とする。定常状態の運動方程式は
\begin{equation}
 \boldsymbol v=-\frac{e\tau}{m}(\boldsymbol E+\boldsymbol v\times\boldsymbol B),\qquad
 v_x=-\frac{e\tau}{m}(E_x+Bv_y),\quad
 v_y=-\frac{e\tau}{m}(E_y-Bv_x).
\end{equation}
電流密度は$\boldsymbol J=-ne\boldsymbol v$なので
\begin{equation}
 \begin{pmatrix}1&x\\-x&1\end{pmatrix}\begin{pmatrix}j_x\\j_y\end{pmatrix}
 =\frac{ne^2\tau}{m}\begin{pmatrix}E_x\\E_y\end{pmatrix}.
\end{equation}
行列を反転すると、問題の形式
\begin{equation}
 \boxed{\widetilde{\sigma}=\frac{\sigma_0}{1+x^2}
 \begin{pmatrix}1&-x\\x&1\end{pmatrix}},\qquad
 \boxed{\sigma_0=\frac{ne^2\tau}{m}}
\end{equation}
を得る。

\section*{[1.2] Hall係数と縦伝導率}
$\boldsymbol j=\widetilde{\sigma}\boldsymbol E$を反転し、$j_y=0$とすると
\begin{equation}
 E_x=\frac{j_x}{\sigma_0},\qquad E_y=-\frac{xj_x}{\sigma_0}.
\end{equation}
従って
\begin{equation}
 \boxed{R_H=-\frac{x}{\sigma_0B}=-\frac{1}{ne}},\qquad
 \boxed{\sigma=\frac{j_x}{E_x}=\sigma_0=\frac{ne^2\tau}{m}}.
\end{equation}
Hall係数が負なのは、担い手の電荷が$-e$であることを反映している。

\section*{[2.1] 電子・正孔の伝導率テンソル}
電子のテンソルは[1.1]の結果であり、正孔は反対符号の電荷を持つのでHall成分の符号が反転する。両者の電流を加えると
\begin{equation}
 \boxed{\widetilde{\sigma}=\frac{e^2\tau/m}{1+x^2}
 \begin{pmatrix}n+p&-(n-p)x\\(n-p)x&n+p\end{pmatrix}}.
\end{equation}
ここで$x=\omega_c\tau$は正の量である。

\section*{[2.2] Hall係数と縦伝導率}
$u=n+p$、$d=n-p$とおく。伝導率行列の逆は
\begin{equation}
 \widetilde{\sigma}^{-1}=\frac{m(1+x^2)}{e^2\tau}
 \frac{1}{u^2+d^2x^2}
 \begin{pmatrix}u&dx\\-dx&u\end{pmatrix}.
\end{equation}
$j_y=0$では$\boldsymbol E=\widetilde{\sigma}^{-1}(j_x,0)^T$なので
\begin{align}
 R_H&=\frac{E_y}{j_xB}
 =-\frac{d}{e}\frac{1+x^2}{u^2+d^2x^2}\\
 &=\boxed{-\frac{n-p}{e}\frac{1+(\omega_c\tau)^2}{(n+p)^2+(n-p)^2(\omega_c\tau)^2}},\\
 \sigma&=\frac{j_x}{E_x}
 =\boxed{\frac{(n+p)e^2\tau}{m}
 \frac{1+[(n-p)/(n+p)]^2x^2}{1+x^2}}.
\end{align}

\section*{[2.3] 低磁場・高磁場Hall係数}
式の$x\to0$および$x\to\infty$の極限は
\begin{equation}
 R_H^{(0)}=-\frac{n-p}{e(n+p)^2},\qquad
 R_H^{(\infty)}=-\frac{1}{e(n-p)}.
\end{equation}
従って
\begin{equation}
 n-p=-\frac{1}{eR_H^{(\infty)}},\qquad
 n+p=\frac{1}{e\sqrt{R_H^{(0)}R_H^{(\infty)}}},
\end{equation}
ここで平方根は正の枝を取った。よって
\begin{equation}
 \boxed{n=\frac{1}{2e}\left[\frac{1}{\sqrt{R_H^{(0)}R_H^{(\infty)}}}-\frac{1}{R_H^{(\infty)}}\right]},\quad
 \boxed{p=\frac{1}{2e}\left[\frac{1}{\sqrt{R_H^{(0)}R_H^{(\infty)}}}+\frac{1}{R_H^{(\infty)}}\right]}.
\end{equation}
非負密度を持つこのモデルの極限値は、同符号であることに加え
\begin{equation}
 0<\frac{R_H^{(0)}}{R_H^{(\infty)}}
 =\frac{d^2}{u^2}\le1
\end{equation}
を満たす。従って復元した$u$と$d$には$u\ge|d|$が成り立ち、$n=(u+d)/2$、$p=(u-d)/2$はともに非負となる。二つのHall係数を入力として使う場合にも、同符号と$|R_H^{(0)}|\le|R_H^{(\infty)}|$の両方が、このモデルとの適合条件である。

\section*{[2.4] 縦伝導率の磁場依存性}
\begin{equation}
 \sigma(0)=\frac{(n+p)e^2\tau}{m},\qquad
 \sigma(\infty)=\frac{(n-p)^2e^2\tau}{m(n+p)}.
\end{equation}
$n,p\ge0$なら$|(n-p)/(n+p)|\le1$なので、$p=0$では$\sigma$は一定、両種が存在し$n\ne p$なら単調に減少して正の高磁場極限へ近づく。代表的な曲線を規格化して示す（混合担体の曲線は$|n-p|/(n+p)=1/2$の例）。
\begin{figure}[htbp]
\centering
\begin{tikzpicture}[xscale=0.95,yscale=0.75,>=Stealth]
 \begin{scope}
  \draw[->] (0,0)--(5.6,0) node[right] {$x=\omega_c\tau$};
  \draw[->] (0,0)--(0,3.0) node[above] {$\sigma/\sigma(0)$};
  \draw (0,1)--(5.1,1);
  \draw (0,0.08)--(0,-0.08) node[below] {$0$};
  \node at (2.7,-0.55) {(a) $p/n=0$};
  \node[left] at (0,1) {$1$};
 \end{scope}
 \begin{scope}[xshift=8.4cm]
  \draw[->] (0,0)--(5.6,0) node[right] {$x=\omega_c\tau$};
  \draw[->] (0,0)--(0,3.0) node[above] {$\sigma/\sigma(0)$};
  \draw[dashed] (0,0.25)--(5.1,0.25);
  \draw[thick,domain=0:5,samples=100] plot (\x,{(1+0.25*\x*\x)/(1+\x*\x)});
  \draw (0,0.08)--(0,-0.08) node[below] {$0$};
  \node at (2.7,-0.55) {(b) $p/n>0,\ p\ne n$};
  \node[left] at (0,1) {$1$};
 \end{scope}
\end{tikzpicture}
\caption{縦伝導率の代表的な磁場依存性。右図では$[(n-p)/(n+p)]^2=1/4$の例を示す。}
\end{figure}

\section*{[2.5] 電子・正孔密度が等しい場合}
$p=n>0$では伝導率テンソルのHall成分が相殺する。[2.1]の行列から$j_y=2ne^2\tau E_y/[m(1+x^2)]$なので、全電流の条件$j_y=0$よりHall電場は$E_y=0$となる。このとき縦伝導率は
\begin{equation}
 \sigma=\frac{2ne^2\tau/m}{1+x^2}\longrightarrow0\qquad(x\to\infty).
\end{equation}
$E_y=0$で各担体の定常運動方程式を解くと、縦速度は
\begin{equation}
 \text{電子：}\quad v_x=-\frac{e\tau E_x}{m(1+x^2)},\qquad
 \text{正孔：}\quad v_x=\frac{e\tau E_x}{m(1+x^2)},
\end{equation}
横速度はともに$v_y=-e\tau xE_x/[m(1+x^2)]$である。横速度は同方向だが電荷の符号が逆なので、等密度では横電流が相殺する。そのため、単一担体の場合のように横のローレンツ力を補償するHall電場が生じない。高磁場では縦速度が$1/(1+x^2)$の因子で抑えられ、電場方向の電流と縦伝導率がゼロへ近づく。
\end{document}
